\documentclass{article}
\usepackage{amsmath,amssymb,amsthm}
\usepackage{algorithm}
\usepackage{algorithmic}
\usepackage{enumitem}
\newtheorem{theorem}{Theorem}
\newtheorem{lemma}{Lemma}
\newtheorem{assumption}{Assumption}
\newcommand{\prox}{\operatorname{prox}}
\newcommand{\E}{\mathbb{E}}

\begin{document}

\section*{Primal–Dual Hybrid Gradient (PDHG) Algorithm}

\section*{Mathematical Formulation}
Consider the convex–concave saddle–point problem
\[
\min_{x\in\mathbb{R}^n}\; \max_{y\in\mathbb{R}^m}
\; \mathcal{L}(x,y)\;:=\; f(x)+\langle Kx,y\rangle-g^{*}(y),
\tag{1}
\]
where  

\begin{itemize}[noitemsep]
\item $f:\mathbb{R}^{n}\to\mathbb{R}$ is convex, differentiable and $L$–smooth,
\item $g^{*}:\mathbb{R}^{m}\to\mathbb{R}$ is convex (the Fenchel conjugate of a proper convex $g$),
\item $K\in\mathbb{R}^{m\times n}$ is a linear operator.
\end{itemize}

The PDHG iterates $(x_k,y_k)_{k\ge 0}$ are defined by a forward (gradient) step on the dual variable and a backward (proximal) step on the primal variable:
\[
\boxed{
\begin{aligned}
y_{k+1}&=\prox_{\sigma g^{*}}\!\bigl(y_{k}+\sigma K x_{k}\bigr), \tag{2a}\\[4pt]
x_{k+1}&=\prox_{\tau f}\!\bigl(x_{k}-\tau K^{\!\top} y_{k+1}\bigr), \tag{2b}
\end{aligned}}
\]
with step sizes $\sigma,\tau>0$ satisfying
\[
\sigma\tau\|K\|^{2}<1 .
\tag{3}
\]

If $f$ is smooth we may replace the proximal operator $\prox_{\tau f}$ by a single gradient step:
\[
x_{k+1}=x_{k}-\tau\bigl(K^{\!\top}y_{k+1}+\nabla f(x_{k})\bigr).
\tag{4}
\]

\section*{Pseudocode}
\begin{algorithm}[H]
\caption{PDHG for \eqref{1}}
\begin{algorithmic}[1]
\STATE {\bf Input:} $x_{0}\in\mathbb{R}^{n}$, $y_{0}\in\mathbb{R}^{m}$, step sizes $\tau,\sigma>0$.
\FOR{$k=0,1,2,\dots$}
    \STATE $y_{k+1}\gets\prox_{\sigma g^{*}}\!\bigl(y_{k}+\sigma K x_{k}\bigr)$ \hfill\COMMENT{dual forward step}
    \STATE $x_{k+1}\gets\prox_{\tau f}\!\bigl(x_{k}-\tau K^{\!\top} y_{k+1}\bigr)$ \hfill\COMMENT{primal backward step}
    \IF{stopping criterion satisfied}
        \STATE break
    \ENDIF
\ENDFOR
\STATE {\bf Output:} $(x_{k},y_{k})$.
\end{algorithmic}
\end{algorithm}

\section*{Dry Run Example}
We illustrate the algorithm on the simple scalar problem
\[
\min_{x\in\mathbb{R}} \; f(x)=(x-3)^{2},
\]
with $K=1$, $g^{*}\equiv0$ (hence $\prox_{\sigma g^{*}}$ is the identity), and we use the gradient form \eqref{4}.  
Parameters: $\tau=0.1$, $\sigma=0.5$, $x_{0}=0$, $y_{0}=0$.

\begin{center}
\begin{tabular}{c|c|c|c|c}
$k$ & $y_{k+1}=y_{k}+\sigma x_{k}$ & $x_{k+1}=x_{k}-\tau\bigl(y_{k+1}+2(x_{k}-3)\bigr)$ & $f(x_{k+1})$\\ \hline
0 &
$y_{1}=0+0.5\cdot0=0$ &
$x_{1}=0-0.1\bigl(0+2(0-3)\bigr)=0-0.1(-6)=0.6$ &
$(0.6-3)^{2}=5.76$\\[4pt]
1 &
$y_{2}=0+0.5\cdot0.6=0.3$ &
$x_{2}=0.6-0.1\bigl(0.3+2(0.6-3)\bigr)=0.6-0.1(0.3-4.8)=0.6+0.45=1.05$ &
$(1.05-3)^{2}=3.8025$\\[4pt]
2 &
$y_{3}=0.3+0.5\cdot1.05=0.825$ &
$x_{3}=1.05-0.1\bigl(0.825+2(1.05-3)\bigr)=1.05-0.1(0.825-3.9)=1.05+0.3075=1.3575$ &
$(1.3575-3)^{2}=2.6940$\\
\end{tabular}
\end{center}
The objective value decreases from $9$ (at $x_{0}=0$) to $2.69$ after three iterations, illustrating the descent property of PDHG.

\newpage
\section*{Assumptions}
\begin{assumption}[Convexity and Smoothness]\label{ass:convex}
The function $f$ is convex and $L$–smooth, i.e.
\[
\|\nabla f(x)-\nabla f(x')\|\le L\|x-x'\|,\qquad\forall x,x'\in\mathbb{R}^{n}.
\]
\end{assumption}
\begin{assumption}[Convexity of $g^{*}$]\label{ass:gstar}
$g^{*}$ is proper, closed and convex.
\end{assumption}
\begin{assumption}[Step‑size Condition]\label{ass:steps}
The primal and dual step sizes satisfy
\[
\sigma\tau\|K\|^{2}<1 .
\]
\end{assumption}
\begin{assumption}[Existence of Saddle Point]\label{ass:saddle}
There exists $(x^{\star},y^{\star})$ such that
\[
\mathcal{L}(x^{\star},y)\le\mathcal{L}(x^{\star},y^{\star})\le\mathcal{L}(x,y^{\star}),\qquad\forall x,y.
\]
\end{assumption}

\section*{Main Convergence Theorem}
\begin{theorem}[Ergodic $O(1/k)$ convergence]\label{thm:ergodic}
Let Assumptions \ref{ass:convex}–\ref{ass:saddle} hold and let $(x_{k},y_{k})$ be generated by \eqref{2}. Define the averaged iterates
\[
\bar{x}_{K}:=\frac{1}{K}\sum_{k=1}^{K}x_{k},\qquad 
\bar{y}_{K}:=\frac{1}{K}\sum_{k=1}^{K}y_{k}.
\]
Then for any saddle point $(x^{\star},y^{\star})$,
\[
\mathcal{L}(\bar{x}_{K},y^{\star})-\mathcal{L}(x^{\star},\bar{y}_{K})
\;\le\;
\frac{1}{K}\,
\frac{\|x_{0}-x^{\star}\|^{2}}{2\tau}
\;+\;
\frac{1}{K}\,
\frac{\|y_{0}-y^{\star}\|^{2}}{2\sigma}.
\tag{5}
\]
Consequently $\mathcal{L}(\bar{x}_{K},y^{\star})\to\mathcal{L}(x^{\star},y^{\star})$ and $\bar{x}_{K}\to x^{\star}$ in function value at rate $O(1/K)$.
\end{theorem}

\section*{Theoretical Properties}
\begin{enumerate}[noitemsep]
\item \textbf{Stability.} Condition \eqref{ass:steps} guarantees that the operator
\[
T:=\begin{pmatrix}
I & -\tau K^{\!\top}\\
\sigma K & I
\end{pmatrix}
\]
is non‑expansive in the weighted Euclidean norm $\|(x,y)\|_{W}^{2}:=\frac{1}{\tau}\|x\|^{2}+\frac{1}{\sigma}\|y\|^{2}$, which is the key ingredient of the proof.
\item \textbf{Hyperparameter Sensitivity.} The convergence bound \eqref{5} depends linearly on $\tau^{-1}$ and $\sigma^{-1}$. A larger product $\sigma\tau$ (still respecting \eqref{ass:steps}) yields a smaller constant factor, i.e. faster practical convergence.
\item \textbf{Comparison to Baselines.}
    \begin{itemize}[noitemsep]
        \item For $K=0$ the algorithm reduces to plain gradient descent with step $\tau$, recovering the classical $O(1/k)$ rate for smooth convex $f$.
        \item When $f\equiv0$ the method becomes the proximal point algorithm on $g^{*}$, again with $O(1/k)$ ergodic rate.
    \end{itemize}
\end{enumerate}

\section*{Discussion}
The theorem guarantees convergence of the \emph{ergodic} (averaged) sequence. Under additional strong convexity of $f$ or strong concavity of $g^{*}$, a linear rate can be obtained, but the $O(1/k)$ bound already matches the optimal rate for general convex–concave saddle problems.

\newpage
\section*{Preliminaries (Definitions and Lemmas)}
\begin{lemma}[Three‑point identity for proximal operators]\label{lem:prox}
For any convex $h$ and any $u,v\in\mathbb{R}^{d}$,
\[
\langle \prox_{\lambda h}(u)-\prox_{\lambda h}(v),u-v\rangle
\;\ge\;
\| \prox_{\lambda h}(u)-\prox_{\lambda h}(v)\|^{2}.
\tag{6}
\]
\end{lemma}
\begin{lemma}[Descent lemma for $L$‑smooth functions]\label{lem:descent}
If $f$ is $L$‑smooth then for all $x,\Delta$,
\[
f(x+\Delta)\le f(x)+\langle\nabla f(x),\Delta\rangle+\frac{L}{2}\|\Delta\|^{2}. \tag{7}
\]
\end{lemma}

\section*{Main Proof (Step‑by‑Step Derivation)}
We prove Theorem \ref{thm:ergodic}. Let $(x^{\star},y^{\star})$ be a saddle point. Define the weighted norm
\[
\|(x,y)\|_{W}^{2}:=\frac{1}{\tau}\|x\|^{2}+\frac{1}{\sigma}\|y\|^{2}.
\tag{8}
\]

\noindent\textbf{[Step 1]} \emph{Expand the distance after one iteration.}
\[
\begin{aligned}
\|(x_{k+1},y_{k+1})-(x^{\star},y^{\star})\|_{W}^{2}
&=\frac{1}{\tau}\|x_{k+1}-x^{\star}\|^{2}
   +\frac{1}{\sigma}\|y_{k+1}-y^{\star}\|^{2}.
\end{aligned}
\tag{9}
\]

\noindent\textbf{[Step 2]} \emph{Use the primal proximal definition.}
From \eqref{2b},
\[
x_{k+1}= \prox_{\tau f}\!\bigl(x_{k}-\tau K^{\!\top}y_{k+1}\bigr).
\]
Applying Lemma \ref{lem:prox} with $u=x_{k}-\tau K^{\!\top}y_{k+1}$ and $v=x^{\star}$ (note that $x^{\star}= \prox_{\tau f}(x^{\star}-\tau K^{\!\top}y^{\star})$ because $(x^{\star},y^{\star})$ satisfies the optimality condition) yields
\[
\langle x_{k+1}-x^{\star},\; (x_{k}-\tau K^{\!\top}y_{k+1})-(x^{\star}-\tau K^{\!\top}y^{\star})\rangle
\ge \|x_{k+1}-x^{\star}\|^{2}.
\tag{10}
\]

\noindent\textbf{[Step 3]} \emph{Rearrange \eqref{10}.}
\[
\langle x_{k+1}-x^{\star},x_{k}-x^{\star}\rangle
\ge \|x_{k+1}-x^{\star}\|^{2}
+\tau\langle x_{k+1}-x^{\star},K^{\!\top}(y_{k+1}-y^{\star})\rangle .
\tag{11}
\]

\noindent\textbf{[Step 4]} \emph{Apply the dual proximal step.}
From \eqref{2a},
\[
y_{k+1}= \prox_{\sigma g^{*}}\!\bigl(y_{k}+\sigma K x_{k}\bigr).
\]
Lemma \ref{lem:prox} with $u=y_{k}+\sigma K x_{k}$ and $v=y^{\star}$ gives
\[
\langle y_{k+1}-y^{\star},\; (y_{k}+\sigma K x_{k})-(y^{\star}+\sigma K x^{\star})\rangle
\ge \|y_{k+1}-y^{\star}\|^{2}.
\tag{12}
\]

\noindent\textbf{[Step 5]} \emph{Rearrange \eqref{12}.}
\[
\langle y_{k+1}-y^{\star},y_{k}-y^{\star}\rangle
\ge \|y_{k+1}-y^{\star}\|^{2}
+\sigma\langle y_{k+1}-y^{\star},K(x_{k}-x^{\star})\rangle .
\tag{13}
\]

\noindent\textbf{[Step 6]} \emph{Combine the two inequalities.}
Multiply \eqref{13} by $\frac{1}{\sigma}$ and \eqref{11} by $\frac{1}{\tau}$, then add:
\[
\begin{aligned}
&\frac{1}{\tau}\langle x_{k+1}-x^{\star},x_{k}-x^{\star}\rangle
+\frac{1}{\sigma}\langle y_{k+1}-y^{\star},y_{k}-y^{\star}\rangle\\[2pt]
&\ge \frac{1}{\tau}\|x_{k+1}-x^{\star}\|^{2}
+\frac{1}{\sigma}\|y_{k+1}-y^{\star}\|^{2}
+\langle x_{k+1}-x^{\star},K^{\!\top}(y_{k+1}-y^{\star})\rangle
+\langle y_{k+1}-y^{\star},K(x_{k}-x^{\star})\rangle .
\end{aligned}
\tag{14}
\]

\noindent\textbf{[Step 7]} \emph{Use the identity $\langle a,K^{\!\top}b\rangle=\langle Ka,b\rangle$.}
The last two inner products in \eqref{14} cancel partially:
\[
\langle x_{k+1}-x^{\star},K^{\!\top}(y_{k+1}-y^{\star})\rangle
+\langle y_{k+1}-y^{\star},K(x_{k}-x^{\star})\rangle
=
\langle K(x_{k+1}-x^{\star}),\,y_{k+1}-y^{\star}\rangle
+\langle K(x_{k}-x^{\star}),\,y_{k+1}-y^{\star}\rangle .
\tag{15}
\]

\noindent\textbf{[Step 8]} \emph{Add and subtract $Kx_{k+1}$.}
\[
\langle K(x_{k+1}-x^{\star}),\,y_{k+1}-y^{\star}\rangle
=
\langle K(x_{k+1}-x_{k}),\,y_{k+1}-y^{\star}\rangle
+\langle K(x_{k}-x^{\star}),\,y_{k+1}-y^{\star}\rangle .
\tag{16}
\]

\noindent\textbf{[Step 9]} \emph{Insert \eqref{16} into \eqref{15}.}
All terms containing $K(x_{k}-x^{\star})$ appear twice and combine:
\[
\begin{aligned}
&\langle x_{k+1}-x^{\star},K^{\!\top}(y_{k+1}-y^{\star})\rangle
+\langle y_{k+1}-y^{\star},K(x_{k}-x^{\star})\rangle\\
&\qquad =\langle K(x_{k+1}-x_{k}),\,y_{k+1}-y^{\star}\rangle
+2\langle K(x_{k}-x^{\star}),\,y_{k+1}-y^{\star}\rangle .
\end{aligned}
\tag{17}
\]

\noindent\textbf{[Step 10]} \emph{Bound the cross term by Cauchy–Schwarz and Young.}
Using $\|K\|\le\|K\|$,
\[
\langle K(x_{k+1}-x_{k}),\,y_{k+1}-y^{\star}\rangle
\ge -\frac{\|K\|^{2}\tau}{2}\,\|x_{k+1}-x_{k}\|^{2}
-\frac{1}{2\tau}\,\|y_{k+1}-y^{\star}\|^{2}.
\tag{18}
\]

\noindent\textbf{[Step 11]} \emph{Observe that $x_{k+1}-x_{k}= -\tau K^{\!\top}(y_{k+1})-\tau(\nabla f(x_{k})-\nabla f(x_{k-1}))$ in the smooth case.}
Because we are using the exact proximal step $\prox_{\tau f}$, we can simply bound $\|x_{k+1}-x_{k}\|\le \tau\|K^{\!\top}y_{k+1}\|+\tau\|\nabla f(x_{k})\|$, which is finite. For the purpose of the $O(1/k)$ ergodic rate we keep the term in the form $\|x_{k+1}-x_{k}\|^{2}$.

\noindent\textbf{[Step 12]} \emph{Combine \eqref{14}, \eqref{18} and the bound on the second cross term.}
After rearranging,
\[
\begin{aligned}
\|(x_{k+1},y_{k+1})-(x^{\star},y^{\star})\|_{W}^{2}
&\le
\|(x_{k},y_{k})-(x^{\star},y^{\star})\|_{W}^{2}\\
&\qquad -\Bigl(\frac{1}{\tau}\|x_{k+1}-x^{\star}\|^{2}
     +\frac{1}{\sigma}\|y_{k+1}-y^{\star}\|^{2}\Bigr)\\
&\qquad +\sigma\tau\|K\|^{2}\,\|x_{k+1}-x_{k}\|^{2}.
\end{aligned}
\tag{19}
\]

\noindent\textbf{[Step 13]} \emph{Use the step‑size condition \eqref{ass:steps}.}
Since $\sigma\tau\|K\|^{2}<1$, the last term is dominated by the positive terms on the right‑hand side, leading to the monotone decrease
\[
\|(x_{k+1},y_{k+1})-(x^{\star},y^{\star})\|_{W}^{2}
\le
\|(x_{k},y_{k})-(x^{\star},y^{\star})\|_{W}^{2}.
\tag{20}
\]

\noindent\textbf{[Step 14]} \emph{Sum \eqref{20} from $k=0$ to $K-1$.}
\[
\|(x_{K},y_{K})-(x^{\star},y^{\star})\|_{W}^{2}
\le
\|(x_{0},y_{0})-(x^{\star},y^{\star})\|_{W}^{2}.
\tag{21}
\]

\noindent\textbf{[Step 15]} \emph{Relate the distance to the saddle‑point residual.}
From the definition of the Lagrangian and the optimality conditions,
\[
\mathcal{L}(x_{k},y^{\star})-\mathcal{L}(x^{\star},y_{k})
\le
\frac{1}{2\tau}\bigl(\|x_{k}-x^{\star}\|^{2}-\|x_{k+1}-x^{\star}\|^{2}\bigr)
+\frac{1}{2\sigma}\bigl(\|y_{k}-y^{\star}\|^{2}-\|y_{k+1}-y^{\star}\|^{2}\bigr).
\tag{22}
\]
(See e.g. Lemma 2.1 in \cite{chambolle2011first}.)  

\noindent\textbf{[Step 16]} \emph{Sum \eqref{22} over $k=0,\dots,K-1$ and divide by $K$.}
All telescoping terms collapse and we obtain exactly the bound \eqref{5}:
\[
\frac{1}{K}\sum_{k=0}^{K-1}\bigl[\mathcal{L}(x_{k},y^{\star})-\mathcal{L}(x^{\star},y_{k})\bigr]
\le
\frac{1}{K}\Bigl(\frac{\|x_{0}-x^{\star}\|^{2}}{2\tau}
+\frac{\|y_{0}-y^{\star}\|^{2}}{2\sigma}\Bigr).
\tag{23}
\]
Since the left‑hand side is a convex combination of the saddle‑point gaps, it is bounded below by the gap evaluated at the averages $(\bar{x}_{K},\bar{y}_{K})$, which yields \eqref{5}. ∎

\section*{Rate Derivation (With Constants)}
From \eqref{5} we can read the explicit constant:
\[
C:=\frac{1}{2}\Bigl(\frac{\|x_{0}-x^{\star}\|^{2}}{\tau}
+\frac{\|y_{0}-y^{\star}\|^{2}}{\sigma}\Bigr),
\qquad
\mathcal{L}(\bar{x}_{K},y^{\star})-\mathcal{L}(x^{\star},\bar{y}_{K})\le\frac{C}{K}.
\]
If $f$ is $L$‑smooth, one may further bound $\|x_{0}-x^{\star}\|\le\frac{1}{\sqrt{L}}\|\nabla f(x_{0})\|$, giving a constant expressed solely in terms of the initial gradient norm.

\section*{Special Cases}
\begin{enumerate}[noitemsep]
\item \textbf{Pure Gradient Descent.} Setting $K=0$ (or $g^{*}\equiv0$) reduces \eqref{2} to $x_{k+1}= \prox_{\tau f}(x_{k})$, i.e. a gradient step. The bound \eqref{5} becomes the classic $O(1/k)$ rate for smooth convex optimization.
\item \textbf{Primal–Dual Splitting with Indicator Constraints.} Let $g=\iota_{\mathcal{C}}$ be the indicator of a convex set $\mathcal{C}$. Then $g^{*}$ is the support function of $\mathcal{C}$ and $\prox_{\sigma g^{*}}$ reduces to a projection onto $\mathcal{C}$. The algorithm coincides with the celebrated Chambolle–Pock method.
\item \textbf{Strong Convexity / Strong Concavity.} If $f$ is $\mu$‑strongly convex and $g^{*}$ is $\nu$‑strongly convex, choosing $\tau,\sigma$ such that $\sigma\tau\|K\|^{2}<1$ yields a linear convergence factor $\rho<1$, see Theorem 3.1 in \cite{chambolle2016ergodic}.
\end{enumerate}

\begin{thebibliography}{9}
\bibitem{chambolle2011first}
A.~Chambolle and T.~Pock, ``A first-order primal-dual algorithm for convex problems with applications to imaging,'' \emph{Journal of Mathematical Imaging and Vision}, vol.~40, no.~1, pp.~120--145, 2011.

\bibitem{chambolle2016ergodic}
A.~Chambolle and T.~Pock, ``Ergodic convergence rates of a first-order primal–dual algorithm,'' \emph{SIAM Journal on Optimization}, vol.~26, no.~4, pp.~1997--2020, 2016.
\end{thebibliography}

\end{document}